\originalsection{18.315：作业5（原卷 Fall 2005）}
本组收录原题3、4. 原题3本身允许明确说明能够完整分类的子类;
下面完整分类最大度不超过2的情形, 再给出更高度数的例族.

\begin{exercise}\label{mit:18315-hw5-q3}
\textbf{18.315, 原卷 Fall 2005, 作业5, 题3.}
尝试分类全部有限、连通、可平面、顶点传递的图.
若不能找出全部, 应明确说明自己能够完整分类的子类.
\end{exercise}
\sourceof{官方 HW5 第1页题3; 编者独立解答, 非官方答案; 图均为有限简单图}
\begin{solution}
图的自同构是保持邻接关系的顶点双射. 顶点传递指任意 $u,v$ 都有一个自同构把
$u$ 映到 $v$. 自同构把邻点集双射到邻点集, 因而所有顶点度数相同.
这先把问题限制在正则图中; 正则性本身还不保证顶点传递.

\textbf{最大度不超过2的完整分类.}
答案恰为 $K_1,K_2$ 和 $C_m$ ($m\ge3$).
若公共度数为0, 连通性给出仅一个顶点. 若为1, 从任一顶点到它的唯一邻点后
不能再走到第三点, 连通性给出 $K_2$.
若为2, 从一个顶点开始, 每次沿着不同于来边的边继续走.
有限性使某个顶点首次重复; 在首次重复之前各点不同.
首次重复必须回到起点, 否则被重复点已有两条此前走过的圈边及一条进入路径的边,
与度数2矛盾. 所得圈的每一点已经用尽两条邻边, 不可能再与圈外顶点相邻.
连通性于是给出整图就是这个圈.
反过来, $K_1,K_2$ 显然可平面且顶点传递;
$C_m$ 画成正 $m$ 边形, 循环旋转把任一点移到任一点, 也满足全部条件.

\textbf{若干更高度数的经典例子.}
$K_4$ 把一顶点放在三角形内部、连向三角形的三个顶点即可平面绘制,
而任意顶点置换都是自同构. 八面体图可记作 $K_{2,2,2}$:
把其中四点画成一个圈, 剩下两点分别放在球面圈的两侧,
各连向圈上四点. 三个二元部之间可以任意置换, 每部内也可以交换两点,
因此该图顶点传递.

还有无穷例族 $C_m\mathbin\square K_2$ ($m\ge3$).
顶点写作 $(i,\varepsilon)$, $i\in\mathbb Z/m\mathbb Z$, $\varepsilon\in\{0,1\}$;
边连接 $(i,\varepsilon)$ 与 $(i+1,\varepsilon)$, 以及 $(i,0)$ 与 $(i,1)$.
把两个圈画成内外同心的 $m$ 边形, 对应顶点径向相连, 即得到无交叉画法.
映射 $(i,\varepsilon)\mapsto(i+t,\varepsilon+s)$,
其中第一坐标模 $m$、第二坐标模2, 都是自同构, 并且传递地作用在全部顶点上.
$m=4$ 给出立方体图. 这些三度、四度例子说明上面的低度分类范围确实有限;
本解答不把这一例表称为所有顶点传递平面图的分类.
\end{solution}

\begin{exercise}\label{mit:18315-hw5-q4}
\textbf{18.315, 原卷 Fall 2005, 作业5, 题4.}
证明或否证: 每张没有分离三角形的平面三角剖分都有 Hamilton 圈.
这里分离三角形指图中的三角形不是一个面边界.
\end{exercise}
\sourceof{官方 HW5 第1页题4; 编者独立解答, 非官方答案}
\begin{solution}
命题成立. 以下在球面上讨论有限简单三角剖分, 至少有三个顶点;
“没有分离三角形”就是说每个三角形至少有一侧不含图的顶点,
且那一侧是一个三角形面. 我们使用正文已说明的 Jordan 圈分离工具,
而不调用任何外部 Hamilton 定理.

\textbf{归纳对象.}
一个简单圈围成的闭盘连同盘中全部顶点、边称为三角盘;
其内面都是三角形. 边界上连接两个不相邻边界顶点、而内部落在盘内的边称为弦.
证明下面的加强命题:

若三角盘的边界可分为两条或三条首尾相接的无弦弧,
则任取两个不同的分界点 $A,B$, 都有一条从 $A$ 到 $B$ 的路,
经过盘中及边界全部顶点各一次.
这里“无弦弧”包括其两个端点: 不准有盘内弦连接该弧的任意两点.
所有盘都取自这张没有分离三角形的球面三角剖分.
对盘的总顶点数归纳, 三个顶点时直接成立.

\textbf{邻扇事实.}
沿一个边界顶点 $v$ 的盘内邻点次序写作 $u_0,u_1,\ldots,u_s$.
连续邻点相邻, 因为它们与 $v$ 围成的面为三角形.
这些邻点依次连接形成一条路; 删除 $v$ 及其邻扇的三角形后,
它代替原来的边界段 $u_0vu_s$.
若 $v$ 除 $u_0,u_s$ 外没有别的盘边界邻点, 新路的中间顶点都是盘内顶点,
剩余区域是一个较小三角盘, 或仅剩一条边.
新路在剩盘中没有弦: 假如非连续的 $u_i,u_j$ 由剩盘内弦相连,
三角形 $vu_iu_j$ 一侧有中间的邻点, 另一侧有被弦隔出的边界点,
成为分离三角形, 矛盾.
邻扇中连续的小三角形本来是面, 所以删除邻扇不会遗漏夹在其中的顶点.

先说明一个供最后一类使用的\textbf{扇区拼接事实}.
设一个盘的边界依次为
\[
A,u_1,\ldots,u_s,x,y,w_{t-1},\ldots,w_1,A,
\]
把 $W=A,w_1,\ldots,w_{t-1},y$ 看成从 $A$ 到 $y$ 的边界弧,
把 $U=A,u_1,\ldots,u_s,x$ 看成另一路; $t\ge1$, $A\ne y$,
$x,y$ 之间是边界边.
假设 $W$ 无弦, 并且 $U$ 上除 $A$ 外每一点都邻接 $W$ 的某一点.
只要归纳命题已对严格较小的盘成立,
便可从 $A$ 走到 $y$, 恰好经过原盘除 $u_1,\ldots,u_s,x$ 外的全部顶点.

为证此事实, 画出全部 $U$ 到 $W$ 的邻接边, 包括最后的 $xy$.
它们的端点次序单调: 若 $u_i$ 早于 $u_j$, 它在 $W$ 上的邻点不可能
晚于 $u_j$ 的邻点, 否则两边端点交替, Jordan 分离迫使交叉.
同一个 $u_i$ 可以有多个 $W$ 邻点, 与下一个 $U$ 顶点的邻点只能按此顺序衔接.
这里把 $U$ 的最后一点 $x$ 一并当作待删的 $U$ 顶点,
并把 $A$ 到第一个 $U$ 点的边及 $xy$ 也标入.
取画出\textbf{全部}这些边后所得的极小扇区, 即其内部不再含一条 $U$--$W$ 边.
每区在 $U$ 一侧只含一个待删点, 或两个相邻待删点:
若含三个, 中间点必有到 $W$ 的边, 它就会进一步细分该区.
更明确地, 一点区由同一个 $U$ 点的两个相继 $W$ 邻点界定;
两点区由前一个 $U$ 点的最后 $W$ 邻点与后一个 $U$ 点的最先 $W$ 邻点界定.
它们在 $W$ 上的区间按从 $A$ 到 $y$ 的次序排列, 端点可以重合.
下面逐一删去区中的 $U$ 点.

只有一个待删顶点 $v$ 时, 扇区边界为 $avb$ 加上 $W$ 的 $a$--$b$ 段,
其中 $a\ne b$. 极小性保证 $v$ 没有该段内部的 $W$ 邻点,
否则相应已画边还会分区. 因而删去 $v$ 的邻扇后,
余盘边界确由一条无弦邻点链与这段 $W$ 组成, 不会裂为多盘.
有两个待删顶点 $v,v'$ 时, 扇区边界为 $avv'b$ 加上 $W$ 的 $a$--$b$ 段.
先设 $a\ne b$. 极小性同样保证 $v$ 在 $W$ 段上只邻接 $a$,
$v'$ 只邻接 $b$, 包括不能另接相对端点 $b,a$.
边 $vv'$ 的扇区内三角形面第三点记为 $z$;
$z$ 必在扇区严格内部: 若在 $W$ 的内部或端点,
$vz,v'z$ 中至少一条就是尚在该区内部的已画跨边, 与极小性矛盾.
删去二者邻扇后, 余边界由 $a$ 到 $z$ 的 $v$ 邻点链,
$z$ 到 $b$ 的 $v'$ 邻点链及 $W$ 的 $a$--$b$ 段组成.
两条邻点链除 $z$ 外不会相遇: 若还有一个共同点 $z'$,
则它同样在扇区内部, 三角形 $vv'z'$ 的扇区侧必包含紧贴 $vv'$ 的面顶点 $z$,
而 $W$ 的 $a$--$b$ 段在另一侧. 因 $a\ne b$, 另一侧有顶点,
这成为分离三角形, 矛盾. 所以二邻扇相接且删除后确实只余一个盘.
两条链也各自无弦, 原因正是上述邻扇事实.
若两点区 $a=b$, 区边界就是三角形 $avv'$.
$W$ 的另一端点位于区外, 因没有分离三角形, 区内没有顶点;
删除后只留下连接点 $a$, 作为长度为零的路, 不再调用盘归纳.
余区域若仅有边 $ab$, 就直接沿它走; 否则它是严格较小的三角盘,
边界有两条或三条无弦弧. 归纳给出经过余盘全部点的 $a$--$b$ 路.
依次拼接这些路. 切区的公共边本身只有一个 $U$ 端点与一个 $W$ 端点,
删除 $U$ 点后, 相邻区只共享相应的 $W$ 连接点;
所有新增邻点链的内部点严格属于各自扇区, 不可能在另一区再次出现.
由单调端点序, 不相邻区若共享同一个 $W$ 点, 中间各区的 $W$ 区间只能全为该点,
即均是上述空三角形区. 跳过这些零长路时保留该连接点一次,
因此拼接路没有重复, 并覆盖全部剩点. 若某区没有面积、只剩一条 $W$ 边,
就沿该边衔接, 同样不把它作为归纳盘.
这证明了扇区拼接事实.

现在证明归纳步骤. 边界按一个方向写成
\[
A,a_1,\ldots,a_\alpha,B,b_1,\ldots,b_\beta,
C,c_1,\ldots,c_\gamma,A.
\]
三条指定弧分别是 $A$--$B$, $B$--$C$, $C$--$A$.
若原来只有两弧, 从 $B$--$A$ 弧的内部任选一点作为 $C$;
因为边界至少三点, 该弧总可选择为有内部点的一弧.
弧上的 $a_i,b_j,c_k$ 列表可以为空.
把 $C$ 也记作 $c_0$, 以便统一写弦端点.
以下写子盘边界时, 首尾同点且没有边的零长弧均省去;
例如 $y=C$ 时的 $C$--$y$ 弧不计作第三弧.
仅剩一条边的情形直接沿边走, 不把它冒充三角盘;
实际应用归纳的盘始终有两条或三条非零无弦弧及不同分界点.

\textbf{第一类: 某个 $a_i$ 邻接 $C$ 或一个 $c_j$.}
若只有 $a_i$ 到 $b_j$ 的弦, 交换 $A,B$ 并反向读边界,
就化为这里的情形, 最后再反向读所得路即可.
选最小的 $i$, 再在该 $i$ 下选最大的 $j$, 记所选弦为 $xy$, $x=a_i,y=c_j$.
它切出含 $A$ 的盘 $D_A$ 与含 $B$ 的盘 $D_B$.
在 $D_A$ 中, $x$ 只有前一个 $a$ 点(或 $A$)与 $y$ 两个边界邻点,
$y$ 只有 $x$ 与后一个 $c$ 点(或 $A$)两个边界邻点;
这是最小 $i$、最大 $j$ 及原来各弧无弦共同保证的.

至少下面两种情况之一成立:
$y$ 不邻接 $a_{i+1},\ldots,a_\alpha,B$;
或者 $x$ 不邻接 $c_0,c_1,\ldots,c_{j-1}$.
后一列表在 $j=0$ 时为空, 于是后一条件自动成立;
它绝不包括已由所选弦相连的 $y=C$.
若二者都不成立, 所得两条弦的端点分别落在 $x$ 后、$y$ 前,
与 $xy$ 相对的次序使它们互相交替, 不可能同时存在.
前一种情况中, 在 $D_A$ 删除 $x$ 的邻扇.
余盘边界的三弧为原来从 $A$ 到 $x$ 前点的弧、替换邻点链及从 $y$ 到 $A$ 的弧;
它们无弦, 归纳得 $A$--$y$ 路, 覆盖 $D_A$ 除 $x$ 外全部点.
若余盘只是一条边, 直接取该边.
$D_B$ 的三弧为 $y,x,a_{i+1},\ldots,B$, $B,\ldots,C$ 及 $C,\ldots,y$,
由本情况及原无弦条件仍然无弦.
归纳得覆盖 $D_B$ 的 $y$--$B$ 路. 两路只共享 $y$, 拼接即得答案.
后一种情况则在 $D_A$ 删除 $y$ 的邻扇, 得覆盖其余点的 $A$--$x$ 路;
$D_B$ 用三弧 $x,\ldots,B$, $B,\ldots,C$, $C,\ldots,y,x$,
归纳给出 $x$--$B$ 路. 这次只共享 $x$, 同样完成.
以上删除与弦分割的盘都严格小于原盘.

\textbf{第二类: 第一类没有发生, 但 $B$ 邻接一个 $c_j$, 或 $A$ 邻接一个 $b_i$.}
只须考虑前者; 后者交换 $A,B$.
选择最大 $j$, 以 $By$ ($y=c_j$) 分盘.
含 $A$ 的一侧中, $B$ 只邻接前一个 $a$ 点(或 $A$)与 $y$ 两个边界点.
删去 $B$ 的邻扇, 原 $A$ 弧、替换链、$y$ 到 $A$ 的弧为两条或三条无弦弧,
归纳得覆盖这一侧除 $B$ 外的 $A$--$y$ 路.
另一侧用三弧 $yB$, $B,\ldots,C$, $C,\ldots,y$,
归纳得 $y$--$B$ 路, 拼接完成. 空三角扇区仍直接用一条边代替.

\textbf{第三类: 整个边界没有弦.}
在 $B$ 到 $A$ 的边界方向取 $A$ 的前邻点 $z\ne B$,
必要时交换两个方向, 确保这样的点存在.
先走 $Az$, 然后删除 $A$ 的邻扇.
新盘从 $z$ 沿替换邻点链到 $A$ 在另一方向的邻点 $w$,
再沿原边界到 $B$, 再从 $B$ 沿原边界回 $z$.
这三条弧无弦, 对较小盘归纳得到 $z$--$B$ Hamilton 路,
在前面添上 $A$ 即可. 若原盘只有面三角形, 直接走 $A,z,B$;
若 $w=B$, 新盘只需要两弧. 因为原边界无弦,
替换链内部不会碰到别的原边界点, 所以这个删除确实留下一个盘.

\textbf{第四类: 只有 $b_i$ 到 $c_j$ 的弦.}
选 $i$ 最小、在该 $i$ 下选 $j$ 最大的一条, 记为 $xy$, $x=b_i,y=c_j$.
其含 $A,B$ 的一侧盘 $D$ 没有任何边界弦:
别的 $b_hc_k$ 弦若在这一侧, 要么 $h<i$, 要么 $h=i,k>j$,
与选择矛盾; 剩下 $h>i,k>j$ 的可能又与 $xy$ 端点交替.
另一侧盘 $D_C$ 的三条无弦弧为边 $yx$, $x,\ldots,C$ 与 $C,\ldots,y$;
归纳给出覆盖 $D_C$ 的 $y$--$x$ 路.

下面把 $D$ 中的点分配给一个 $A$--$y$ 路和一个 $x$--$B$ 路.
令 $W$ 是 $y,c_{j+1},\ldots,c_\gamma,A$ 这条边界链.
其盘内邻点带是连通的: 每条 $W$ 边的盘内三角形面有一个第三顶点;
相邻两条 $W$ 边的第三顶点, 沿公共边界点的盘内邻扇相连.
从 $x$ 沿 $y$ 的邻扇进入此带, 再沿这些相连邻扇,
可以到达 $A$ 的一个盘内邻点.
由于 $D$ 没有边界弦, 带中的这些点除起点 $x$ 外都在盘内部.
因此在“$x$、所有邻接 $W$ 的盘内点、$A$”诱导的图中,
取一条最短的 $x$--$A$ 路.
除其最后一个盘内点外, 此路没有点邻接 $A$, 否则可直接接 $A$ 而缩短.

令 $z$ 为 $A$ 在朝 $B$ 方向的边界邻点, 即 $a_1$ 或 $B$.
把上述路的最后一条到 $A$ 的边改为沿 $A$ 的邻点链走到 $z$,
取路的朝 $B$ 一侧的那段邻扇. 得到一条简单的 $x$--$z$ 路 $P$:
新增点全邻接 $A$, 不会等于此前不邻接 $A$ 的点,
而且它们除 $z$ 外仍为盘内点.
此路的每个内部点都邻接 $W$, 第一个邻接 $A$ 的点之后的全部点都邻接 $A$.

若 $P$ 上两点之间还有位于其 $B$ 侧的弦, 用该弦替换所夹的路段.
每次严格减少点数, 直到得到一条在 $B$ 侧无弦的 $x$--$z$ 路 $Q$.
上述邻接 $W$ 与尾段邻接 $A$ 的性质都保留, 因为只删除点而不新增点.
由 $Q$, $z,\ldots,B$ 与 $B,\ldots,x$ 围出的盘 $D_B$ 有三条无弦弧,
且缺少 $A$, 所以归纳得到覆盖它的 $x$--$B$ 路.

设 $q$ 是 $Q$ 上从 $x$ 起第一个邻接 $A$ 的点.
$Q$ 的 $q$--$z$ 尾段及 $A$ 围成的邻扇里没有其他顶点:
其中每个连续小三角形 $Aq_rq_{r+1}$ 都为面,
否则它有尾段外的边界点在另一侧, 成为分离三角形.
因而尚未归入 $D_B$ 的全部点恰在由
$A,q,\ldots,x,y,c_{j+1},\ldots,c_\gamma,A$ 围出的盘中,
并且要删除 $q,\ldots,x$, 因为它们已在 $D_B$ 或 $D_C$ 的路上.
对这个盘应用扇区拼接事实: 其中 $U=A,q,\ldots,x$ 的每一点除 $A$ 外
均邻接 $W'=A,c_\gamma,\ldots,c_{j+1},y$, 最后一项 $x$ 邻接 $y$;
$W'$ 无弦. 得到恰覆盖其余点的 $A$--$y$ 路.
依次接上 $D_C$ 的 $y$--$x$ 路及 $D_B$ 的 $x$--$B$ 路,
三部分只在接点相遇, 无重复且覆盖原盘所有顶点.
扇区拼接中每个盘都删去了至少一个 $U$ 点, 其余归纳盘也严格小于原盘,
所以不存在循环调用. 第四类及归纳命题均证毕.

最后任选原球面三角剖分的一个面 $ABC$, 取其另一侧作为三角盘.
边界的三条单边弧自动无弦, 加强命题给出覆盖整图的 $A$--$B$ 路.
再加边 $BA$ 就是 Hamilton 圈. 三个顶点时直接是 $K_3$;
以上证明同时包括 $K_4$, 并没有暗中假设顶点数大于4.

原论文定位为 H. Whitney, A theorem on graphs,
Annals of Mathematics 32(2), 1931, 378--390,
\href{https://doi.org/10.2307/1968197}{原文链接}.
这里的盘归纳与邻扇构造使用本书重新叙述的定义与证明;
不附受限论文扫描件, 不以论文引用代替上述推导.
\end{solution}
